%% ma: L12-B18-0020
%% lop: 12
%% chuong: 6
%% bai: 18
%% dang: 12.18.2
%% loai: DS
%% muc: ['TH', 'TH', 'TH', 'TH']
%% dap_an: "SSSS"
%% nguon: "(NBV)-10 ĐỀ KIỂM TRA HK2-TOÁN 12-DUNG CHUNG 3 SÁCH.pdf (H03, Đề số 3), Phần 2 Câu 2"
%% kien_thuc: "Bài 18 (xác suất có điều kiện, biến cố đối)"
%% trang_thai: chinh-thuc
%% ghi_chu: "Đáp án gốc khớp (SSSS). Bỏ ảnh lớp học (ảnh trang trí)."
\begin{ex}
Một lớp học có $17$ học sinh nam và $24$ học sinh nữ. Cô giáo gọi ngẫu nhiên lần lượt $2$ học sinh (có thứ tự) lên trả lời câu hỏi. Xét các biến cố:

$A$: ``Lần thứ nhất cô giáo gọi $1$ học sinh nam'';

$B$: ``Lần thứ hai cô giáo gọi $1$ học sinh nữ''.
\choiceTF
{$P(B\mid A)=0{,}575$.}
{$P\left(B\mid\overline A\right)=0{,}6$.}
{$P\left(\overline B\mid A\right)=0{,}425$.}
{$P\left(\overline B\mid\overline A\right)=0{,}4$.}
\loigiai{Lớp có $17+24=41$ học sinh.

Nếu lần đầu gọi học sinh nam ($A$) thì còn $40$ học sinh, gồm $16$ nam và $24$ nữ: $P(B\mid A)=\dfrac{24}{40}=0{,}6$ và $P\left(\overline B\mid A\right)=\dfrac{16}{40}=0{,}4$.

Nếu lần đầu gọi học sinh nữ ($\overline A$) thì còn $40$ học sinh, gồm $17$ nam và $23$ nữ: $P\left(B\mid\overline A\right)=\dfrac{23}{40}=0{,}575$ và $P\left(\overline B\mid\overline A\right)=\dfrac{17}{40}=0{,}425$.

Vậy cả bốn mệnh đề đều sai.}
\end{ex}
